\documentclass[../main.tex]{subfiles}
\begin{document}
\chapter{Eigenvectors and Eigenvalues}
\section{Polynomials}
\begin{definition}[Polynomial expression] A \textit{polynomial} over a field $\F$ is an \textbf{expression}: \[
    f(t) = a_n t^n + \cdots + a_1 t + a_0
  \]
  for some $n \geq 0$ where $a_i \in \F$ are called the \textit{coefficients} of $f$.
\end{definition}
\begin{remark}
  At this point, we are \textbf{not} thinking of polynomials as maps from $\F \to \F$, we are just thinking of them as finite sequences of coefficients $a_0, \ldots, a_n$ in $\F$ of arbitrary length.
\end{remark}
\begin{remark}[Notation]
  We denote the set of polynomial expressions over $\F$ as $\F[t]$.
\end{remark}
\begin{definition}[Degree]
  The degree of $f(t) \in \F[t]$, denoted $\deg(f)$, is the largest $m$ such that $a_m \neq 0$.
  We call this $a_m$ the \textit{leading coefficient of $f$}.
\end{definition}
\begin{remark}
  By convention, we define the degree of the zero polynomial $f(t) = 0$ to be $-\infty$, i.e. $\deg(0) = -\infty$.
\end{remark}
\begin{definition}[Monic]
  We say that a polynomial $f(t) \in \F[t]$ is \textit{monic} if its leading coefficient is 1.
\end{definition}
We can define arithmetic on $\F[t]$.
For $f, g \in \F[t]$, we can write:
\[
  f(t) = \sum_{i = 0}^{n} a_i t^{i} \text{ and } g(t) = \sum_{j = 0}^{m} b_j t^{j}
\]
and we then define $f + g$ and $f \cdot g$ as:
\begin{align*}
  (f + g)(t) &\text{ has coefficient } a_k + b_k \text{ of $t^k$}\\
  \text{ and } (f \cdot g)(t) &\text{ has coefficient } \sum_{i = 0}^{k} a_i b_{k - i} \text{ of $t^k$}
\end{align*}
\begin{remark}[Remarks]
  \begin{itemize}
    \item From the definitions, we see that both $+$ and $\cdot$ are commutative operations on $\F[t]$.
    \item We see that:
      \[
        \deg(f + g) \leq \max\{\deg(f), \deg(g)\}
      \]
      as the leading coefficients of the polynomials could cancel.
      But, in the case of $\cdot$:
      \[
        \deg(f \cdot g) = \deg(f) + \deg(g)
      \]
      as the leading coefficients of $f$ and $g$ are non-zero.
  \end{itemize}
\end{remark}
\begin{definition}[Divisible]
  For $f, g \in \F[t]$, we say that $f$ \textit{divides} $g$, written $f \mid g$, if there exists $h \in \F[t]$ such that $g = f \cdot h$.
\end{definition}
For $\lambda \in \F$, we \textit{evaluate} a polynomial $f \in \F[t]$ using the expression:
\[
  f(\lambda) = \sum_{i = 0}^{n} a_i \lambda^i \in \F
\]
Evaluation of a polynomial also passes through the operations $+$ and $\cdot$, i.e. $(f + g)(\lambda) = f(\lambda) + g(\lambda)$ and $f(f \cdot g)(\lambda) = f(\lambda)g(\lambda)$.

\begin{remark}[Warning]
  This distinction between a polynomial as a finite collection of coefficients and as something we can evaluate at a particular $\lambda \in \F$ is important.

  If $\F$ is finite, then $\F[t]$ is \textbf{not} isomorphic to the space of polynomial functions $\F \to \F$.
  We see that $\F[t]$ is infinite dimensional as the set $\{t^n: n \in \N_0\}$ is a basis for $\F[t]$.
  However the set of all functions from $\F \to \F$ is finite dimensional as we can think of each function as a column vector in $\F^{|\F|}$.
  Thus as the space of polynomial functions $\F \to \F$ is a subspace of this, it must also be finite dimensional.
\end{remark}
\begin{proposition}[Euclidian algorithm for polynomials]
  For all $f, g \in \F[t]$ and $g \neq 0$, there exists $q, r \in \F[t]$ such that: \label{EAP}
  \[
    f = q \cdot g + r
  \]
  with $\deg(r) < \deg(g)$
\end{proposition}
\begin{proof}
  Let $n = \deg(f)$ and $m = \deg(g)$ and write:
  \[
    f(t) = \sum_{i = 0}^{n} a_i t^{i} \text{ and } g(t) = \sum_{j = 0}^{m} b_j t^{j}
  \]

  If $m > n$, we take $q = 0$ and $r = f$ as then $\deg(r) = \deg(f) < \deg(g)$.
  Otherwise, if $n \geq m$, define $\widetilde{f}(t) = f(t) - (\frac{a_n}{b_m}) t^{n - m} g(t)$.
  The leading coefficient of $t^n$ is $a_n - \frac{a_n}{b_m} \cdot b_m = 0$, and so since $\deg(f) \geq \deg(g)$, there cannot be any terms of greater degree, thus $\deg(f') < \deg(f)$.
  We then repeat the above inductively on $f'$ and $g$ until we reach the case $m > n$.
\end{proof}
\begin{corollary}[Bezout's Lemma]
  If $f_1, \ldots, f_n \in \F[t]$ have no common factor of degree $\geq 1$, then there exists $p_1, \ldots, p_k \in \F[t]$ such that:
  \[
    \sum_{i = 1}^{k} f_i(t)p_i(t) = 1
  \]
\end{corollary}
\begin{proof}
  Proved in the same was as for integers using the Euclidean algorithm in IA Numbers and Sets.
\end{proof}
\begin{lemma}
  For any $\lambda \in \F$, $f(\lambda) = 0$ if and only if $(t - \lambda) \mid f(t)$. \label{polyFac}
\end{lemma}
\begin{proof}
  \begin{proofdirection}{Assume $f(\lambda) = 0$}
    Setting $g(t) = t - \lambda$, we apply \cref{EAP} to $f$ and $g$.
    Thus there exists $q, r \in \F[t]$ such that:
    \[
      f = q \cdot g + r
    \]
    with $\deg(r) < \deg(g) = 1$.
    Thus $r = a$ for some $a \in \F$.
    Evaluating both sides at $\lambda$, we see that:
    \[
      0 = f(\lambda) = q(\lambda) \cdot g(\lambda) + a = 0 + a \implies a = 0
    \]
    So $f = q \cdot g$, i.e. $t - \lambda = g \mid f$.
  \end{proofdirection}
  \begin{proofdirection}{Assume $(t - \lambda) \mid f$}
    So there exists $h \in \F[t]$ such that $f = (t - \lambda) \cdot h$.
    Evaluating at $\lambda$, we have $f(\lambda) = (\lambda - \lambda) h(\lambda) = 0$.
  \end{proofdirection}
\end{proof}
\begin{definition}[Root]
  We say that $\lambda \in \F$ is a \textit{root} of \textit{multiplicity} at least $e \in \N_0$ if:
  \[
    (t - \lambda)^{e} \mid f
  \]
  That is, the multiplicity of a root $\lambda$ is the greatest $e \in \N$ for which this holds.
\end{definition}
\begin{corollary}
  If $\deg(f) = n \geq 0$, then $f$ has at most $n$ roots counted with multiplicity. \label{rootCount}
\end{corollary}
\begin{proof}
  We induct on the degree of $f$, $n = \deg(f)$.
  \quietinduction
  {$n = 0$}{
    So $f(t) = a$ for some $a \in \F \setminus \{0\}$, thus $f$ has zero roots.
  }
  {$n - 1$}{}
  {$n$}{
    Suppose $f$ has a root $\mu \in \F$, then by \cref{polyFac}, $(t - \mu) \mid f$.
    So we can write:
    \[
      f(t) = (t - \mu) \cdot g(t)
    \]
    for some $g(t) \in \F[t]$.

    We see that $n = \deg(f) = \deg(t - \mu) + \deg(g) \implies \deg(g) = n - 1$.
    So by hypothesis, $g$ has at most $n - 1$ roots counted with multiplicity.
    We then observe that $\lambda \in \F$ is a root of $f$ if and only if it is $\mu$ or it is a root of $g$.
    Thus $f$ has at most $n - 1 + 1 = n$ roots counted with multiplicity.
  }
\end{proof}
\begin{corollary}
  If $\deg(f), \deg(g) < n$ and there exists distinct $\lambda_1, \ldots, \lambda_n \in \F$ such that:
  \[
    f(\lambda_i) = g(\lambda_i)\ \forall i \in \{1, \ldots, n\}
  \]
  then $f = g$.
\end{corollary}
\begin{proof}
  So $f - g$ has $n$ roots, but $\deg(f - g) \leq \max\{\deg(f), \deg(g)\} < n$.
  If $\deg(f - g) \geq 0$, then this would contradict \cref{rootCount}.
  Thus $\deg(f - g) = -\infty$, i.e. $f - g = 0 \implies f = g$.
\end{proof}
\begin{theorem}[Fundemental Theorem of Algebra]
  If $f \in \C[t]$ of degree $\deg n \geq 1$, then $f$ has exactly $n$ roots in $\C$ counted with multiplicity.
\end{theorem}
\begin{proof}
  See IB Complex Analysis
\end{proof}
\end{document}
